\begin{proposition}[有限时间谱核的归一化与黄金律时间窗]
把所有末态的失谐写为 $\delta$（单位 $\si{\per\second}$），并把耦合平方与末态测度合并成非负密度
$J(\delta)$，使一阶总跃迁概率为
\begin{equation}
P^{(2)}(T)=\int_{\mathbb R}J(\delta)T^2
\sinc^2\!\left(\frac{\delta T}{2}\right)\,\dd\delta.
\label{eq:feqo-finite-time-spectral-density-probability}
\end{equation}
若 $J$ 有界且在 $\delta=0$ 连续，则
$P^{(2)}(T)/T\to2\pi J(0)$。
这个极限只是二阶表达式的谱极限；用于真实动力学还要求该时间窗内初态未显著耗尽。
\end{proposition}
\begin{proof}
令 $K_T(\delta)=T\sinc^2(\delta T/2)/(2\pi)$。
它非负，且由 $\int_{\mathbb R}(\sin x/x)^2\dd x=\pi$ 得
$\int K_T\dd\delta=1$。对任意 $\Omega>0$，
\[
\int_{|\delta|>\Omega}K_T(\delta)\dd\delta
\le\frac{4}{\pi T\Omega},
\]
因为 $\sin^2\le1$。若 $\omega_J(\Omega)
=\sup_{|\delta|\le\Omega}|J(\delta)-J(0)|$，则
\[
\left|\frac{P^{(2)}(T)}T-2\pi J(0)\right|
\le2\pi\omega_J(\Omega)+\frac{16\|J\|_\infty}{T\Omega}.
\]
先固定足够小的 $\Omega$，再令 $T\Omega$ 足够大，即得到收敛。
这里没有把离散末态变成连续谱：若 $J$ 含孤立 delta 质量，以上连续性条件不成立。
\end{proof}

\begin{exbox}[从短时二次增长到恒定速率：一个可积谱核]
设电子投影后的末态失谐谱在共振附近可用
\[
J(\delta)=\frac{g_*^2a}{\pi(\delta^2+a^2)},\qquad a>0,
\]
表示，其中 $g_*,a$ 都是角频率尺度。失谐变量可正可负；该洛伦兹密度是工作窗模型，
不是把环境的绝对频率延伸到负值。它的相关核由 Fourier 积分给出
$C(\tau)=\int J(\delta)\ee^{\ii\delta\tau}\dd\delta
=g_*^2\ee^{-a|\tau|}$。
将有限时间振幅平方写成两个时间积分，得到
\begin{align*}
P^{(2)}(T)&=\int_0^T\dd t\int_0^T\dd s\,C(t-s)\\
&=2g_*^2\int_0^T(T-\tau)\ee^{-a\tau}\dd\tau\\
&=\frac{2g_*^2}{a^2}(aT-1+\ee^{-aT}).
\end{align*}
于是
\begin{equation}
P^{(2)}(T)=
\begin{cases}
g_*^2T^2[1-aT/3+O((aT)^2)],&aT\ll1,\\
\Gamma_*T-\Gamma_*/a+O(\ee^{-aT}),&aT\gg1,
\end{cases}
\qquad \Gamma_*=\frac{2g_*^2}{a}.
\label{eq:feqo-lorentzian-spectrum-golden-window}
\end{equation}
黄金律需要 $a^{-1}\ll T\ll\Gamma_*^{-1}$，所以存在宽时间窗还要求 $g_*\ll a$。
例如取 $g_*/a=0.05$、$aT=20$，得到 $P^{(2)}\simeq0.095$，
黄金律给 $\Gamma_*T=0.100$；两者相差 $0.005$，正是有限建立时间造成的修正。
若窄共振使 $g_*\gtrsim a$，上述两侧时间尺度不再分离，便应保留环境模的联合动力学。
\end{exbox}
